\documentclass[10pt,a4paper]{amsart}
\usepackage[margin=19mm]{geometry}
\usepackage{amsmath,amssymb,mathtools}
\usepackage[scr=rsfs,cal=euler]{mathalfa}
\usepackage{xfrac}
\usepackage[unicode,hidelinks,bookmarks=false]{hyperref}
\newcommand{\ie}{i.e.\ }
\newcommand{\cf}{cf.\ }
\newcommand{\resp}{respectively\ }
\newcommand{\Subsection}{\subsection}
\providecommand{\nobreakdash}{\nobreak\hbox{-}}
\setlength{\emergencystretch}{2em}
\multlinegap0pt
\usepackage{bm}
\usepackage{bbm}
\usepackage{dsfont}
\usepackage{xspace}
\usepackage{cancel}
\usepackage{mathtools}
\usepackage{amsthm}
\makeatletter
\@ifundefined{thm}{\newtheorem{thm}{Thm}}{}
\@ifundefined{prop}{\newtheorem{prop}[thm]{Proposition}}{}
\makeatother
\DeclareMathOperator{\bal}{\mathrm{bal}}
\newcommand{\Q}{\mathbf{Q}}
\newcommand{\Sel}{Sel}
\newcommand{\Vsdag}{{\mathbf{V}_{Q_0}^\dagger}}
\newcommand{\cc}{\mathbf{c}}
\newcommand{\ch}{\xi}
\newcommand{\ord}{{{\rm ord},{\rm ord}}}
\newcommand{\ppbar}{\bar{\mathfrak{p}}}
\newcommand{\pp}{\mathfrak{p}}
\newcommand{\rH}{H}
\newcommand{\relstr}{{{\rm rel},{\rm str}}}
\newcommand{\unb}{{\boldsymbol{f}}}
\title[Diagonal cycles and anticyclotomic Iwasawa theory of modular]{Diagonal cycles and anticyclotomic Iwasawa theory of modular forms}
\author{Francesc Castella, Kim Tuan Do}
\date{Source: arXiv:2303.06751v2 (CC BY 4.0)}
\begin{document}
\maketitle

\noindent\emph{Source and licence.} Diagonal cycles and anticyclotomic Iwasawa theory of modular forms. Version arXiv:2303.06751v2 (CC BY 4.0), distributed under the Creative Commons Attribution 4.0 International licence, as recorded in the arXiv licence metadata for this exact version and shown on the abstract page https://arxiv.org/abs/2303.06751v2. Full rights notice: licenses/arxiv-2303.06751-CC-BY-4.0.txt.\par\smallskip

\noindent\textit{Editorial scope.} Counted unit: the Prop whose source label is \texttt{prop:factor-S} in the article (source0.tex lines 3445--3454, source label \texttt{prop:factor-S}), delivered together with the article's own complete original proof of it (source0.tex lines 3456--3504, one \texttt{proof} environment of 49 lines). No mathematics is written, repaired or re-derived: statement and proof are the article's own text line for line. Required context retained uncounted: eq:dec-V-first (lines 3422--3427); eq:shapiro (lines 3439--3441). Every definition, display and label of those lines is retained unchanged. Omitted from this package and not redistributed: 1--3421, 3428--3438, 3442--3444, 3455--3455, 3505--4658. Nothing inside a retained excerpt is omitted or altered: every retained body is delivered exactly as the article's lines are written, so no omission receipt of any kind accompanies this package. Citations: the retained excerpts carry no citation command at all, so no bibliography entry of the article is transcribed and no reference is redistributed. Reference adaptation: none was needed and none was made; every reference inside the delivered unit resolves inside it, so no unresolved reference or citation is produced. Printed slips kept literal: no printed word, symbol, hypothesis or number of the counted statement or of its proof was altered. Layout and class: the article's own class is \texttt{amsart} with packages including inputenc, geometry, amssymb, latexsym, amsmath, amsthm, amscd, graphicx, dsfont, stmaryrd, bm, comment, hyperref, xy; the wrapper is the lane's pinned \texttt{amsart} editorial wrapper at 10pt on A4, which restates the theorem-like environments the retained text needs and adds the article's own macro definitions that the retained text needs (\texttt{\textbackslash Q}, \texttt{\textbackslash Sel}, \texttt{\textbackslash Vsdag}, \texttt{\textbackslash bal}, \texttt{\textbackslash cc}, \texttt{\textbackslash ch}, \texttt{\textbackslash ord}, \texttt{\textbackslash pp}, \texttt{\textbackslash ppbar}, \texttt{\textbackslash rH}, \texttt{\textbackslash relstr}, \texttt{\textbackslash unb}), with extra wrapper packages (bm, bbm, dsfont, xspace, cancel, mathtools). The delivered numbering is the unit's own. Assets: no retained span contains a figure environment, a graphics-inclusion command, a PDF-page inclusion, a table or a TikZ picture; no image file is copied into this package, the package declares no asset dependency and no image file is redistributed with it. Licence: version arXiv:2303.06751v2 of this article is distributed under the Creative Commons Attribution 4.0 International licence, as recorded in the arXiv licence metadata for this exact version and shown on the abstract page https://arxiv.org/abs/2303.06751v2. A full rights notice is delivered at licenses/arxiv-2303.06751-CC-BY-4.0.txt. Characters: every retained excerpt is pure ASCII, so the delivered engine, XeLaTeX with its default Latin Modern text font, reports no missing character.
\begin{equation}\label{eq:dec-V-first}
\begin{aligned}
\Vsdag&\simeq T_f^{\vee}\otimes({\mathrm{Ind}}_K^\Q\ch_1^{-1}\psi_0\Psi_{Z_1})\otimes({\mathrm{Ind}}_K^\Q\ch_2^{-1}\psi_0\Psi_{Z_2})\otimes\varepsilon_{\mathrm{cyc}}^{1-r}(\psi_0^{-1}\Psi_{Z_1}^{-1/2}\Psi_{Z_2}^{-1/2}\circ\mathscr{V})\\
&\simeq\bigl(T_f^{\vee}(1-r)\otimes{\mathrm{Ind}}_{K}^\Q(\ch_1^{-1}\ch_2^{-1}\psi_0^{1-\cc}\Psi_{V_1}^{1-\cc})\bigr)\oplus\bigl(T_f^{\vee}(1-r)\otimes{\mathrm{Ind}}_K^\Q(\ch_1^{-1}\ch_2^{-\cc}\Psi_{W_2}^{1-\cc})\bigr),
\end{aligned}
\end{equation}

\begin{equation}\label{eq:shapiro}
\rH^1(\Q,\Vsdag)\simeq\rH^1(K,T_f^{\vee}(1-r)\otimes\ch_1^{-1}\ch_2^{-1}\Psi_{W_1}^{1-\cc})\oplus\rH^1(K,T_f^{\vee}(1-r)\otimes\ch_1^{-1}\ch_2^{-\cc}\Psi_{W_2}^{1-\cc})
\end{equation}

\begin{prop}\label{prop:factor-S}
Under (\ref{eq:shapiro}), the balanced Selmer group $\Sel^{\bal}(\Q,\Vsdag)$ decomposes as
\begin{align*}
\Sel^{\bal}(\Q,\Vsdag)\simeq\Sel_{\relstr}(K,&T_f^{\vee}(1-r)\otimes\ch_1^{-1}\ch_2^{-1}\Psi_{W_1}^{1-\cc})\\ &\oplus\Sel_{\ord}(K,T_f^{\vee}(1-r)\otimes\ch_1^{-1}\ch_2^{-\cc}\Psi_{W_2}^{1-\cc}),
\end{align*}
and the $\unb$-unbalanced Selmer group $\Sel^{\unb}(\Q,\Vsdag)$ decomposes as
\begin{align*}
\Sel^{\unb}(\Q,\Vsdag)\simeq
\Sel_{\ord}(K,&T_f^{\vee}(1-r)\otimes\ch_1^{-1}\ch_2^{-1}\Psi_{W_1}^{1-\cc})\\ &\oplus\Sel_{\ord}(K,T_f^{\vee}(1-r)\otimes\ch_1^{-1}\ch_2^{-\cc}\Psi_{W_2}^{1-\cc}).\end{align*}
\end{prop}

\begin{proof}
From (\ref{eq:dec-V-first}) we see that the balanced local condition is given by
\begin{equation}\label{eq:bal-def}
\begin{aligned}
\mathscr{F}^{\bal}_p(\Vsdag)
\simeq\bigl(T_f^{\vee}(1-r)\otimes\ch_1^{-1}\ch_2^{-1}\Psi_{W_1}^{1-\cc}\bigr)&\oplus\bigl(T_f^{\vee,+}(1-r)\otimes\ch_1^{-1}\ch_2^{-\cc}\Psi_{W_2}^{1-\cc}\bigr)\\
&\oplus\bigl(T_f^{\vee,+}(1-r)\otimes\ch_1^{-\cc}\ch_2^{-1}\Psi_{W_2}^{\cc-1}\bigr).
\end{aligned}
\end{equation}






Put \[\widetilde{\mathbf{V}}_{Q_0}^\dagger=\bigl(T_f^{\vee}(1-r)\otimes\ch_1^{-1}\ch_2^{-1}\Psi_{W_1}^{1-\cc}\bigr)\oplus\bigl(T_f^{\vee}(1-r)\otimes\ch_1^{-1}\ch_2^{-\cc}\Psi_{W_2}^{1-\cc}\bigr),\]
so by (\ref{eq:dec-V-first}) we have
\begin{equation}\label{eq:shapiro-bis}
\rH^1(\Q,\mathbf{V}_{Q_0}^\dagger)\simeq\rH^1(K,\widetilde{\mathbf{V}}_{Q_0}^\dagger).
\end{equation}
Then from (\ref{eq:bal-def}) we obtain
\begin{equation}\label{eq:bal-shapiro}
\begin{split}
\mathscr{F}^{\bal}_\pp(\widetilde{\mathbf{V}}_{Q_0}^\dagger)&\simeq\bigl(T_f^{\vee}(1-r)\otimes\ch_1^{-1}\ch_2^{-1}\Psi_{W_1}^{1-\cc}\bigr)\\
\mathscr{F}^{\bal}_{\ppbar}(\widetilde{\mathbf{V}}_{Q_0}^\dagger)&\simeq\{0\}
\end{split}
\begin{split}
&\oplus\bigl(T_f^{\vee,+}(1-r)\otimes\ch_1^{-1}\ch_2^{-\cc}\Psi_{W_2}^{1-\cc}\bigr),\\
&\oplus \bigl( T_f^{\vee,+}(1-r)\otimes\ch_1^{-1}\ch_2^{-\cc}\Psi_{W_2}^{1-\cc}\bigr),\nonumber
\end{split}
\end{equation}
%
%Thus denoting by $\Sel^{\{p\}}(K,\widetilde{\mathbf{V}}_{Q_0}^\dagger)$ the submodule of $\rH^1(K,\widetilde{\mathbf{V}}_{Q_0}^\dagger)$ consisting of classes unramified outside $p$, we deduce that under the isomorphism (\ref{eq:shapiro-bis}) the balanced Selmer group $\Sel^{\bal}(\Q,\mathbf{V}_{Q_0}^\dagger)$ corresponds to the kernel of the restriction map
%\begin{align*}
%\Sel^{\{p\}}(K,\widetilde{\mathbf{V}}_{Q_0}^\dagger)
%&\rightarrow\frac{\rH^1(K_\pp,\widetilde{\mathbf{V}}_{Q_0}^\dagger)}{\rH^1(K_\pp,(T_f^{\vee}(1-r)\otimes\ch_1^{-1}\ch_2^{-1}\Psi_{W_1}^{1-\cc})\oplus(T_f^{\vee,+}(1-r)\otimes\ch_1^{-1}\ch_2^{-\cc}\Psi_{W_2}^{1-\cc}))}\\
%&\,\times\frac{\rH^1(K_{\ppbar},\widetilde{\mathbf{V}}_{Q_0}^\dagger)}{\rH^1(K_{\ppbar},T_f^{\vee,+}(1-r)\otimes\ch_1^{-1}\ch_2^{-\cc}\Psi_{W_2}^{1-\cc})},
%\end{align*}
which yields the claimed description of $\Sel^{\bal}(\Q,\Vsdag)$. 

Similarly, we find that the $\unb$-unbalanced local condition is given by
\begin{equation}
\begin{aligned}
\mathscr{F}^{\unb}_\pp(\widetilde{\mathbf{V}}_{Q_0}^\dagger)&\simeq\bigl(T_f^{\vee,+}(1-r)\otimes\ch_1^{-1}\ch_2^{-1}\Psi_{W_1}^{1-\cc}\bigr)\oplus\bigl(T_f^{\vee,+}(1-r)\otimes\ch_1^{-1}\ch_2^{-\cc}\Psi_{W_2}^{1-\cc}\bigr),\\
\mathscr{F}^{\unb}_{\ppbar}(\widetilde{\mathbf{V}}_{Q_0}^\dagger)&\simeq\bigl(T_f^{\vee,+}(1-r)\otimes\ch_1^{-1}\ch_2^{-1}\Psi_{W_1}^{1-\cc}\bigr)\oplus\bigl(T_f^{\vee,+}(1-r)\otimes\ch_1^{-1}\ch_2^{-\cc}\Psi_{W_2}^{1-\cc}\bigr),\nonumber
\end{aligned}
\end{equation}
from where the claimed description of $\Sel^\unb(\Q,\Vsdag)$ follows.
\end{proof}

\end{document}
